\scanpage{58}
\mspage{VI}
% archive foliation: 57
\deleted{\unsure{On est}}
Récurrence sur $\dim X$.
\inserted{trivial si $\dim\leq0$.}
\deleted{\unsure{Utilisant}}
\deleted{\unsure{on a}}
lemme 3, on est ramené au cas où $X$ est lisse et
\inserted{connexe. Si $\dim X=1$,}
\deleted{\unsure{admet un morphisme lisse surjectif}}
c'est facile (réduction au cas compact par lemme de Abhyankar)
\deleted{\unsure{et}}
c'est dû à Serre-Lang.
\deleted{\unsure{D'une manière}}

Si $\dim X\geq2$, on est ramené \deleted{\unsure{au cas}}
\inserted{\unsure{par}}
lemme 3, au cas où il y a une fibration \inserted{lisse} $X\to Y$, $\dim Y<\dim X$, \deleted{\unsure{et}}
\unsure{qui est} à fibres géom. connexes non vides, de dim 1.
On applique alors le lemme 1 et la hyp. de récurrence.

\margin{?}
\uline{Prop} Soit $k$ un corps. On suppose que pour tt \deleted{\unsure{$X$}}
courbe \inserted{\unsure{alg.}} lisse $X$ sur $k$, le nb de rev. étales de $X$ de degré donné est fini.
Alors la même chose est vraie pour \uline{tt} $X$ de t.f. sur $k$.
